% ==================== 第 1 章 绪论 ====================
\chapter{绪论}
\label{ch:intro}

\section{研究背景与动机}

随着信息系统对数据机密性要求的不断提高，密码学已成为现代软件系统的基础设施。然而，密码工程的实践经验反复表明：系统的安全性往往并不取决于所选密码算法的强度，而取决于密钥能否在其整个生命周期中得到妥善管理。换言之，密码算法相对成熟，真正的工程难点在于密钥的生成、存储、使用、轮换与销毁。美国国家标准与技术研究院（NIST）在其密钥管理建议中明确指出，密钥管理是密码系统中最复杂且最易出错的部分之一\cite{nist80057}。

密钥管理系统（Key Management System, KMS）正是为应对这一难题而生。它集中管理密钥的全生命周期，向上层应用提供加解密、密钥轮换等服务，并通过访问控制与审计保证密钥使用的合规性。在云计算与多租户场景中，KMS 已成为各大云服务商的核心安全组件。然而，KMS 自身的安全性高度依赖其运行环境：当承载 KMS 的主机（Host）不可信时——例如主机操作系统被攻破、运维人员越权访问，或进程内存遭到窥探——若主密钥（Master Key）与业务 API 处于同一进程，攻击者一旦控制该进程，便可读出主密钥，进而解密其保护的全部数据加密密钥（DEK），造成系统性的数据泄露。

为缩小这一暴露面，可信执行环境（Trusted Execution Environment, TEE）提供了一条可行路径。TEE 通过硬件机制在不可信的主机之上划出一块“小范围可信计算基（Trusted Computing Base, TCB）”，使其中的代码与数据即便面对拥有更高权限的攻击者也能保持机密性与完整性\cite{sabt2015tee}。Intel SGX、ARM TrustZone 等技术正是这一思想的代表\cite{costan2016sgx, pinto2019trustzone}。将密钥的敏感操作放入 TEE，可在 Host 失陷时仍保护主密钥不被窃取。本文正是沿此思路，探索在 Host 不可信前提下的轻量级 KMS 设计与实现。

\section{问题定义}

本文研究的核心问题可表述为：

\begin{quote}
\itshape 在 Host 不可信的前提下，如何设计并实现一套可运行的密钥管理系统，使主密钥与数据密钥的敏感操作不暴露在 Host 的内存与环境中？
\end{quote}

该问题包含三个相互关联的子问题：其一，如何组织密钥层次，使主密钥的暴露面最小化；其二，如何将敏感密码操作隔离至可信组件，并保证 Host 与可信组件之间通信的机密性、完整性与抗重放；其三，如何在保证上述安全目标的同时，提供完整的密钥生命周期管理、访问控制与审计能力。

需要特别说明的是，受限于本科毕业设计的资源与目标，本文不依赖真实硬件 TEE，而是以一个独立进程对 TEE 的关键能力（隔离执行、远程认证、安全信道）进行\textbf{软件模拟}。因此本文在追求设计完整性的同时，将始终明确区分“硬件 TEE 所能提供的保证”与“软件模拟所能提供的保证”，避免对原型能力的过度表述。

\section{研究目标}

围绕上述问题，本文设定如下研究目标：

\begin{enumerate}
  \item 设计“主密钥 $\rightarrow$ 数据密钥 $\rightarrow$ 业务数据”的两层密钥结构，并以信封加密实现。
  \item 实现密钥的全生命周期管理，包括创建、轮换、禁用、吊销与销毁。
  \item 将主密钥的持有与数据加解密等敏感操作迁移至模拟 Enclave，使其不出现在 Host 中。
  \item 建立 Host 与 Enclave 之间的远程认证与安全通信协议。
  \item 提供 JWT 鉴权、基于角色的访问控制与审计日志，并给出可复现的原型与自动化测试。
\end{enumerate}

\section{主要贡献}

本文的主要贡献概括如表~\ref{tab:contrib} 所示。

\begin{table}[H]
\centering
\caption{本文主要贡献}
\label{tab:contrib}
\small
\begin{tabularx}{\textwidth}{cY}
\toprule
编号 & 贡献 \\
\midrule
1 & 设计并实现了一套轻量级 KMS 架构，集成 RBAC、审计与版本化密钥轮换。 \\
2 & 提出并实现了 Host/Enclave 分离的信封加密方案，使明文数据密钥不出 Enclave。 \\
3 & 实现了简化的远程认证与基于 AES-GCM 的会话信道，并支持序列号防重放。 \\
4 & 建立了系统的威胁模型与安全需求，并诚实界定了软件模拟 TEE 的能力边界。 \\
5 & 提供了可一键启动的原型系统与八项自动化集成测试，保证结果可复现。 \\
\bottomrule
\end{tabularx}
\end{table}

\section{论文组织结构}

本文共分八章，组织如下：第~\ref{ch:intro}~章为绪论，介绍研究背景、问题与目标；第~\ref{ch:background}~章综述相关技术与理论基础，包括认证加密、信封加密、KMS、TEE 与访问控制；第~\ref{ch:threat}~章建立系统的威胁模型与安全需求；第~\ref{ch:design}~章给出系统总体设计，包括架构、密钥层次、模块划分与数据模型；第~\ref{ch:protocol}~章详述关键协议与算法实现，包括远程认证握手与加解密全链路时序；第~\ref{ch:impl}~章介绍系统实现与部署；第~\ref{ch:eval}~章给出测试与实验评估；第~\ref{ch:conclusion}~章总结全文并展望未来工作。
